\section{Methodology}\label{sec:methodology}
We consider a generative model of the form
\begin{equation}\label{eq:model}
  p_\theta(x, z) = p_\theta(x \mid z)\, p(z),
\end{equation}
where \(x \in \mathcal{X}\) denotes observed variables, \(z \in \mathcal{Z}\) latent variables, and \(\theta \in \Theta\) the model parameters.

\subsection{Structural Identifiability}
A central question is whether \(\theta\) can be recovered uniquely from the marginal distribution \(p_\theta(x)\).
We say that the model \eqref{eq:model} is \emph{conditionally structurally identifiable} if, for any \(\theta, \theta' \in \Theta\),
\begin{equation*}
  p_\theta(x) = p_{\theta'}(x) \;\Longrightarrow\; \theta \sim \theta',
\end{equation*}
where \(\sim\) denotes equivalence up to a permitted transformation group (e.g., permutations of latent components).

\subsection{Variational Estimator}
Direct maximum-likelihood estimation is intractable for the model class we consider.
We therefore introduce a variational posterior \(q_\phi(z \mid x)\) and optimize the standard evidence lower bound,
\begin{equation}\label{eq:elbo}
  \mathcal{L}(\theta, \phi) = \mathrm{E}_{q_\phi(z \mid x)}\!\left[\log p_\theta(x \mid z)\right] - \mathrm{KL}\!\left(q_\phi(z \mid x) \,\Vert\, p(z)\right).
\end{equation}
Under mild regularity conditions, gradient ascent on \eqref{eq:elbo} converges to a stationary point of the marginal likelihood at rate \(\mathcal{O}(1/\sqrt{T})\) after \(T\) iterations.

\paragraph{Implementation.}
The encoder \(q_\phi\) and decoder \(p_\theta(x \mid z)\) are parameterized as neural networks; we use the reparameterization trick to obtain low-variance gradient estimates.
